\scanpage{517}
\mspage{428}
% B01
\deleted{b) Si $\pi_0,\pi_1,\pi$ sont choisis comme indiqué, il faut voir que $i_s(E_{s\pi})\subset E_\pi$, i.e que tt $x\in\mathrm{Ob}(E_{s\pi})$ est dans $E_\pi$. Il faut donc prouver que pour tt système inductif $(y_i)_{i\in I}$ grd devant $\pi$, \inserted{de limite $y$ dans $E$, (qui existe par hypothèse sur $\pi_1$ et $\pi\geq\pi_1$),} la flèche canonique
{\csname tagsleft@true\endcsname
\[\varinjlim_I\mathrm{Hom}_E(x,y_i)\longrightarrow\mathrm{Hom}_E(x,y)\tag*{\unsure{(**)|($w$)}}\]}
est bijective.}

% B02
b) Soit $(x_i)_{i\in I}$ un système inductif grd devant $\pi_1$, et $(s_i)_{\unsure{j|i}\in I}$ l'image de ce système dans $C$. \deleted{On verra plus \unsure{b|bas}} Nous prouvons plus bas le

% B03
\uline{Lemme} 4.15.3. Soit $C$ une catégorie, \inserted{$\pi$ un cardinal} tels que $\mathrm{card}\,C\leq\pi$, et que $C$ soit sans projecteurs.

% B04
% slanted
\margin{Cf démonstration page 452.}
% B05
Alors pour tout système inductif $(s_i)_{i\in I}$ grd devant $\pi$ dans $C$, existe une partie cofinale \deleted{de} $J$ de $I$ telle que \deleted{les morphismes de \unsure{de} transition du} le système induit $(s_i)_{i\in J}$ soit \uline{constant}, i.e. tel que les $s_i$ soient égaux à un $s\in C$, et que les morphismes de transition soient \deleted{\unsure{l'identité de}} \inserted{tous} \unsure{égaux|identiques} à $\mathrm{id}_s$.

% B06
% slanted
\margin{Sous \unsure{réserve|réserves} de la validité du \unsure{corollaire|coroll.} de 4.8.13. Mais c'est à présent établi, sous une \unsure{forme|hyp.} \unsure{suffisante|plus suffisante|moins restrictive}, \unsure{cf.} \unsure{4.15.6 a) 1°]|4.15.6 a) 1°|4.15.6 a)}}

% B07
Mais Attention, on s'appuie pour prouver 4.15.3 sur \inserted{le th.} 4.8.13 (pages 212--216), que je n'ai tjrs pas prouvé ! Dans le cas où $C$ est rigide (i.e. $\mathrm{End}_C(s)=\{\mathrm{id}_s\}\quad\forall s\in\mathrm{Ob}\,\unsure{E|C}$), la démonstration est cependant immédiate, cf plus bas.

% archive sheet number: 516
